%=====================================================================
%  CE-1115 Seguridad de la Informacion -- Proyecto, Parte 1
%  Entregable 3: Matriz de control de acceso
%  Cliente: Verdeazul Carbono S.A.
%  Tecnologico de Costa Rica -- Escuela de Ingenieria en Computadores
%  Prof. MSc. Andres Vargas Rivera -- Semestre II, 2026
%
%  Compilar dos veces:  pdflatex P1_Matriz_grupo6.tex
%  Las respuestas del cliente estan en la minuta (Anexo A de P1_Politicas).
%=====================================================================
\documentclass[11pt,letterpaper]{article}

\usepackage[utf8]{inputenc}
\usepackage[T1]{fontenc}
\usepackage[spanish,es-noshorthands,es-tabla]{babel}
\usepackage{mathptmx}
\renewcommand{\ttdefault}{ptm}
\usepackage[letterpaper,margin=2.4cm,headheight=15pt]{geometry}
\usepackage{microtype}
\usepackage{amsmath,amssymb}
\usepackage{xcolor}
\usepackage{colortbl}
\usepackage{booktabs,tabularx,array,longtable,multirow}
\usepackage{enumitem}
\usepackage{titlesec}
\usepackage{fancyhdr}
\usepackage{pdflscape}
\usepackage[most]{tcolorbox}
\usepackage{caption}
\usepackage[hidelinks]{hyperref}

%--------------------------- datos del grupo -------------------------
\newcommand{\numgrupo}{6}
\newcommand{\titulodoc}{Matriz de control de acceso}
\newcommand{\titulocorto}{Matriz de control de acceso}
\newcommand{\numentregable}{3}
\newcommand{\fechaentrega}{15 de octubre de 2026}

%--------------------------- estilo ----------------------------------


\titleformat{\section}{\Large\rmfamily\bfseries }{\thesection}{0.6em}{}
\titleformat{\subsection}{\large\rmfamily\bfseries }{\thesubsection}{0.6em}{}
\titlespacing*{\section}{0pt}{1.4em}{0.5em}
\titlespacing*{\subsection}{0pt}{1.0em}{0.3em}
\setlength{\parindent}{0pt}
\setlength{\parskip}{0.55em}
\setlist{leftmargin=1.5em,itemsep=2pt,topsep=3pt,parsep=0pt}
\captionsetup{font=small,labelfont=bf,skip=5pt}

\pagestyle{fancy}
\fancyhf{}
\fancyhead[L]{\small\rmfamily CE-1115 $\cdot$ Proyecto Parte 1 $\cdot$ Grupo \numgrupo}
\fancyhead[R]{\small\rmfamily \titulocorto}
\fancyfoot[C]{\small\rmfamily \thepage}
\renewcommand{\headrulewidth}{0.4pt}
\renewcommand{\headrule}{\hbox to\headwidth{\leaders\hrule height \headrulewidth\hfill}}

\newcolumntype{L}[1]{>{\raggedright\arraybackslash}p{#1}}
\newcolumntype{Y}{>{\raggedright\arraybackslash}X}
\newcommand{\cab}[1]{\textbf{\rmfamily #1}}

% Trazabilidad con la consulta al cliente (minuta: Anexo A de P1_Politicas)
%   \resp{P-07}: la regla sale de una respuesta del cliente
%   \supu{P-13}: supuesto que el equipo fijo y documento
\newcommand{\resp}[1]{\textrm{\footnotesize\bfseries [minuta #1]}}
\newcommand{\supu}[1]{\textrm{\footnotesize [supuesto #1]}}
\newtcolorbox{decision}{enhanced,breakable,colback=white,colframe=black,
  boxrule=0pt,leftrule=3pt,arc=0pt,left=8pt,right=8pt,top=5pt,bottom=5pt}

% Notacion de la matriz
\newcommand{\cnd}[1]{\textsuperscript{\rmfamily\bfseries #1}}   % condicion
\newcommand{\no}{---}                  % sin acceso
\newcommand{\rol}[1]{\textrm{\bfseries #1}}
\newcommand{\est}[1]{\textsc{\small #1}}

\begin{document}

%=====================================================================
%  PORTADA
%=====================================================================
\begin{titlepage}
\centering
\vspace*{0.6cm}
{\large\scshape Instituto Tecnológico de Costa Rica}\\[0.2cm]
{\large\scshape Escuela de Ingeniería en Computadores}\\[0.2cm]
{\large CE-1115 Seguridad de la Información}\\[1.5cm]

{\rule{\textwidth}{1.4pt}}\\[0.6cm]
{\rmfamily\large Proyecto $\cdot$ Parte 1 $\cdot$ Diseño del control de acceso}\\[0.45cm]
{\rmfamily\bfseries\fontsize{24}{29}\selectfont \titulodoc}\\[0.45cm]
{\rmfamily\large Entregable \numentregable\ de 5}\\[0.6cm]
{\rule{\textwidth}{1.4pt}}\\[1.2cm]

{\large\bfseries Cliente}\\[0.25cm]
{\large Verdeazul Carbono S.A.}\\[0.1cm]
{Certificación de captura de carbono $\cdot$ Pérez Zeledón, Costa Rica}\\[1.2cm]

{\large\bfseries Grupo \numgrupo}\\[0.4cm]
\begin{tabular}{@{}l@{\hspace{1.4cm}}l@{}}
\toprule
\textbf{Integrante} & \textbf{Carné} \\
\midrule
Randall Bryan Bolaños López & 2019043784 \\
Javier Tenorio Cervantes    & 2020065308 \\
Julio Varela Venegas        & 2019008041 \\
Axel Flores Lara            & 2021453573 \\
\bottomrule
\end{tabular}\\[1.2cm]

{\large Prof. MSc. Andrés Vargas Rivera}\\[0.15cm]
{\large Semestre II, 2026}

\vfill
{\large\bfseries Fecha de entrega: \fechaentrega}
\end{titlepage}

\pagenumbering{roman}
\tableofcontents
\newpage
\pagenumbering{arabic}

%=====================================================================
\section{Cómo se lee esta matriz}
%=====================================================================

Esta es la versión comprobable del diseño: cada decisión del documento de control de acceso
(Entregable~2) termina aquí convertida en una celda. Las filas son los
\textbf{sujetos}, las columnas son los \textbf{objetos} y cada celda dice qué puede hacer ese sujeto
sobre ese objeto.

Tres reglas de lectura valen para toda la matriz:

\begin{enumerate}
\item \textbf{Lo que no está escrito está denegado.} El sistema niega por omisión. Una celda con
      --- significa que el sujeto no tiene ningún permiso sobre el objeto, y
      dentro de una celda con permisos, toda operación que no aparezca también se deniega.
\item \textbf{Una condición solo puede quitar.} Un permiso con una marca \cnd{C$n$} se concede
      únicamente si la regla C$n$ de la sección~\ref{sec:reglas} se cumple en ese momento. Las
      condiciones nunca amplían lo que el rol permite.
\item \textbf{La marca \cnd{H} indica vigilancia por historial.} Las lecturas marcadas pasan además por los
      controles H1 a H3 (frecuencia, horario y combinación de ambos), que responden de forma graduada:
      avisan, piden una confirmación adicional o restringen por un tiempo.
\end{enumerate}

Las marcas como \resp{P-07} indican que la regla sale de una respuesta del cliente, recogida en la
minuta (Anexo~A del informe de políticas). Las marcas como \supu{P-13} indican un supuesto que el
equipo fijó y documentó porque el cliente no dio el dato. El resumen está en la sección~\ref{sec:minuta}.

\subsection{Sujetos}

\begin{center}
\small
\begin{tabularx}{\textwidth}{@{}L{1.1cm} L{3.7cm} Y@{}}
\toprule
\cab{Sigla} & \cab{Sujeto} & \cab{Quién es en Verdeazul} \\
\midrule
\rol{TEC} & Técnico de campo & Visita la finca, mide y registra. Trabaja con las fincas que tiene asignadas. \\
\rol{VER} & Verificador & Revisa la medición desde la oficina y la aprueba o la devuelve. Por excepción puede medir, con una habilitación temporal de técnico. \resp{P-01} \\
\rol{SUP} & Supervisor de campo & Coordina al personal de campo: asigna y reasigna las fincas y atiende las alertas de sus técnicos. \resp{P-05} \\
\rol{GER} & Gerencia de Certificación & Responsable del dato. Incluye a la coordinación de certificación: decide los períodos y da la autorización final de las enmiendas y la segunda firma. \resp{P-05, P-07} \\
\rol{ADM} & Administrador de catálogo y cuentas & Personal administrativo: mantiene el catálogo de fincas, registra los períodos y da de alta y de baja las cuentas. \\
\rol{AUD} & Auditor externo & Revisa una vez al año, en sitio o a distancia. Ve todo, exporta evidencia y no modifica nada. \resp{P-18} \\
\rol{DUE} & Dueño de finca & El titular, persona física o sociedad, o su representante autorizado. Consulta lo de sus fincas. \resp{P-15} \\
\rol{SIS} & Procesos del sistema & Emisión de certificados, escritura de la bitácora y motor de alertas. No hay una persona detrás. \\
\rol{OPE} & Operador de plataforma & Personal que la empresa designa para administrar servidores y respaldos. No tiene cuenta en la aplicación. \resp{P-23} \\
\rol{PUB} & Público sin autenticar & Cualquiera que tenga en la mano el código de un certificado, por ejemplo un comprador. \resp{P-04} \\
\bottomrule
\end{tabularx}
\end{center}

Los siete primeros son los roles del modelo RBAC. \rol{SIS} es una identidad de servicio. \rol{OPE} y
\rol{PUB} aparecen en la matriz para dejar escrito qué no pueden hacer: quien custodia la
infraestructura no tiene permisos sobre los datos del negocio, y quien no se identifica solo puede
comprobar un certificado.

\subsection{Objetos}

\begin{center}
\small
\begin{tabularx}{\textwidth}{@{}L{4.3cm} Y@{}}
\toprule
\cab{Objeto} & \cab{Qué contiene} \\
\midrule
Fincas y parcelas & Catálogo: finca, parcelas con su geometría y área, ubicación y titular asociado. \\
Datos personales del dueño & Identificación y contacto del titular y de sus representantes. \\
Períodos de medición & Fechas de inicio, de fin de medición y de cierre de registro, y estado: \est{planificado}, \est{abierto}, \est{cerrado}. \\
Asignaciones & Qué técnico mide qué finca: \emph{ordinaria}, por el período, o \emph{temporal} (apoyo, sustitución o visita extraordinaria), con fecha de fin. \\
Mediciones & Registro de campo por parcela y período, con todas sus versiones, su fecha de medición y de registro, y su estado: \est{borrador}, \est{enviada}, \est{devuelta}, \est{aprobada}, \est{excluida}. \\
Enmiendas & Solicitud, corrección propuesta, revisión técnica y autorización final de un cambio sobre una medición aprobada o de un período cerrado. \\
Certificados & Certificados de carbono emitidos, \est{vigente} o \est{revocado}, con el enlace entre el revocado y su reemplazo. \\
Cuentas y roles & Usuarios, su rol, su vigencia, sus credenciales y los vínculos que declaran con titulares. \\
Bitácora & Registro de quién hizo o consultó qué, cuándo y con qué resultado. \\
Alertas de historial & Avisos que generan los controles H1 a H3, con su nivel y su resolución. \\
\bottomrule
\end{tabularx}
\end{center}

\subsection{Permisos}

Se usan los verbos del dominio en lugar de siglas, para que la matriz se pueda leer proyectada:
\textbf{leer}, \textbf{crear}, \textbf{modificar}, \textbf{eliminar} y \textbf{aprobar}, más los que el negocio
necesita: \textbf{enviar} y \textbf{devolver} una medición; \textbf{solicitar}, \textbf{corregir}, \textbf{revisar} y
\textbf{autorizar} una enmienda; \textbf{emitir} y \textbf{revocar} un certificado; \textbf{abrir} y \textbf{cerrar} un
período; \textbf{proponer} un cambio que otra persona debe aprobar; \textbf{desactivar} una cuenta;
\textbf{resolver} una alerta; \textbf{exportar} evidencia; \textbf{comprobar} un certificado y \textbf{anexar} a la
bitácora.

El permiso \textbf{eliminar} no aparece en ninguna celda de mediciones, enmiendas, certificados,
cuentas ni bitácora. Es una decisión de diseño: en esos objetos nada se borra, todo cambio queda como
una versión nueva o como un cambio de estado.

%=====================================================================
\begin{landscape}
\section{La matriz}
\label{sec:matriz}

\vspace{-0.5em}
{\small\rmfamily
\renewcommand{\arraystretch}{1.1}
\setlength{\tabcolsep}{3.5pt}
\begin{tabularx}{\linewidth}{@{}>{\bfseries}L{1.15cm}|Y|Y|Y|Y|Y|Y|Y|Y|Y|Y@{}}
{\normalfont\scriptsize\rmfamily Sujeto\newline$\downarrow$} & \cab{Fincas y parcelas} & \cab{Datos pers. del dueño} & \cab{Períodos de medición} & \cab{Asigna\-ciones} & \cab{Mediciones} & \cab{Enmiendas} & \cab{Certificados} & \cab{Cuentas y roles} & \cab{Bitácora} & \cab{Alertas de historial} \\
\hline
TEC &
leer\cnd{C1} &
leer\cnd{C1}\newline{\scriptsize solo contacto} &
leer &
leer\cnd{C1}\newline{\scriptsize las propias} &
leer\cnd{C1}\newline crear\cnd{C2}\newline modificar\cnd{C2}\newline enviar\cnd{C2} &
leer\cnd{C1}\newline solicitar\cnd{C7}\newline corregir\cnd{C9} &
\no &
modificar\cnd{C13} &
\no &
\no \\
\hline
VER &
leer\cnd{H} &
\no &
leer &
\no &
leer\cnd{H}\newline aprobar\cnd{C3}\newline devolver\cnd{C3} &
leer\newline solicitar\cnd{C7}\newline revisar\cnd{C10} &
leer\cnd{H} &
modificar\cnd{C13} &
\no &
\no \\
\hline
SUP &
leer\cnd{H} &
\no &
leer &
leer\newline crear\cnd{C19}\newline modificar\cnd{C19} &
leer\cnd{H}\newline{\scriptsize solo estado y fechas} &
leer\newline{\scriptsize solo estado} &
\no &
modificar\cnd{C13} &
\no &
leer\cnd{C15}\newline resolver\cnd{C15} \\
\hline
GER &
leer\cnd{H}\newline aprobar\cnd{C11} &
leer\cnd{H} &
leer\newline abrir\cnd{C14}\newline cerrar\cnd{C14} &
leer &
leer\cnd{H} &
leer\newline solicitar\cnd{C7}\newline autorizar\cnd{C8} &
leer\cnd{H} &
leer\newline aprobar\cnd{C12}\newline modificar\cnd{C13} &
leer &
leer\newline resolver\cnd{C15} \\
\hline
ADM &
leer\cnd{H}\newline proponer\cnd{C11} &
leer\cnd{H}\newline crear\newline modificar\newline eliminar\cnd{C16} &
leer\newline crear\newline proponer\cnd{C14} &
\no &
\no &
\no &
\no &
leer\newline crear\cnd{C12}\newline desactivar\cnd{C12}\newline proponer\cnd{C12}\newline modificar\cnd{C13} &
\no &
\no \\
\hline
AUD &
leer\cnd{C6}\newline exportar\cnd{C18} &
leer\cnd{C6} &
leer\cnd{C6}\newline exportar\cnd{C18} &
leer\cnd{C6}\newline exportar\cnd{C18} &
leer\cnd{C6}\newline exportar\cnd{C18} &
leer\cnd{C6}\newline exportar\cnd{C18} &
leer\cnd{C6}\newline exportar\cnd{C18} &
leer\cnd{C6}\newline modificar\cnd{C13} &
leer\cnd{C6}\newline exportar\cnd{C18} &
leer\cnd{C6}\newline exportar\cnd{C18} \\
\hline
DUE &
leer\cnd{C5} &
leer\cnd{C5} &
leer &
leer\cnd{C5} &
leer\cnd{C5}\newline{\scriptsize solo aprobadas} &
leer\cnd{C5} &
leer\cnd{C5} &
modificar\cnd{C13} &
\no &
\no \\
\hline
SIS &
leer &
leer\newline{\scriptsize solo nombre del titular} &
leer &
leer &
leer &
leer &
emitir\cnd{C4}\newline revocar\cnd{C4} &
\no &
anexar &
crear \\
\hline
OPE & \no & \no & \no & \no & \no & \no & \no & \no & \no & \no \\
\hline
PUB & \no & \no & \no & \no & \no & \no & comprobar\cnd{C20} & \no & \no & \no \\
\end{tabularx}
}

\vspace{0.2em}
{\footnotesize\rmfamily
\textbf{Leyenda.}\quad
--- sin acceso: se deniega.\quad
\cnd{C$n$} permiso sujeto a la regla C$n$ (sección~\ref{sec:reglas}).\quad
\cnd{H} lectura vigilada por los controles de historial H1 a H3.\quad
Toda operación no escrita en una celda está denegada.\quad
\textbf{C17} (abstención por vínculo) condiciona además todo verbo que no sea \emph{leer} de TEC, VER, SUP, GER y ADM
sobre fincas, asignaciones, mediciones y enmiendas.\quad
\emph{Aprobar}, \emph{devolver}, \emph{revisar}, \emph{autorizar}, \emph{abrir}, \emph{cerrar}, \emph{resolver} y \emph{exportar}
exigen reautenticación con segundo factor en los cinco minutos previos.\quad
\textbf{OPE:} su acceso a la infraestructura se trata en el documento de arquitectura.}
\end{landscape}

%=====================================================================
\section{Reglas de las celdas condicionadas}
\label{sec:reglas}
%=====================================================================

Cada regla se escribe de forma que se pueda comprobar con un caso de prueba: dada una situación
concreta, el sistema concede o deniega. En las reglas, $s$ es el sujeto que hace la petición, $m$ una
medición, $e$ una enmienda y $\mathit{Autores}(m)$ es el conjunto de personas que registraron o
modificaron cualquier versión de $m$. La columna \emph{Depende de} indica si la regla mira el
\textbf{estado} del sistema, el \textbf{historial} del sujeto o un \textbf{atributo} del objeto.

{\small
\renewcommand{\arraystretch}{1.25}
\begin{longtable}{@{}L{0.8cm} L{2.5cm} L{8.0cm} L{1.7cm} L{1.95cm}@{}}
\toprule
\cab{Id.} & \cab{Nombre} & \cab{Regla: el permiso se concede solo si\ldots} & \cab{Depende de} & \cab{Celdas} \\
\midrule
\endfirsthead
\toprule
\cab{Id.} & \cab{Nombre} & \cab{Regla: el permiso se concede solo si\ldots} & \cab{Depende de} & \cab{Celdas} \\
\midrule
\endhead
\bottomrule
\endfoot

C1 & Asignación vigente &
existe una asignación de $s$ a la finca del objeto y la fecha actual cae dentro de su vigencia: la
\emph{ordinaria} vale mientras el período esté \est{abierto}; la \emph{temporal}, hasta su fecha de fin.
También si $s$ es el técnico que debe corregir una enmienda en trámite sobre esa finca. Además,
siempre puede leer las mediciones que él mismo registró. \resp{P-11} &
Estado, atributo y tiempo &
TEC: fincas, datos, asignaciones, mediciones, enmiendas \\

C2 & Edición solo en borrador &
se cumple C1, el período de $m$ está \est{abierto}, $m$ está en \est{borrador} o \est{devuelta}, y
$s$ es su autor o el técnico a quien SUP reasignó la finca; en ese caso el cambio es una versión
nueva, con $s$ como autor. La fecha de medición que $s$ declara cae dentro del período y de su
asignación, y no es posterior a la de registro, que pone el servidor. \emph{Enviar} pasa $m$ a
\est{enviada} y la deja fuera del alcance del técnico. \resp{P-22} &
Estado e historial &
TEC: mediciones \\

C3 & Quien registró no aprueba &
$m$ está \est{enviada} y $s \notin \mathit{Autores}(m)$. Vale con el período \est{abierto} o ya
\est{cerrado}: el cierre detiene el registro, no la revisión de lo que se envió a tiempo. La regla
mira todas las versiones: vale aunque $s$ haya cambiado de rol o haya medido con una habilitación
temporal. \resp{P-01, P-06} &
Historial y estado &
VER: mediciones \\

C4 & Emisión única &
la medición está \est{aprobada}, quien la aprobó no está en $\mathit{Autores}(m)$ y \textbf{no existe otro
certificado \est{vigente} para la misma parcela y el mismo período}. \emph{Revocar} solo ocurre por
una enmienda autorizada (C8) que cambia las toneladas certificadas: en una sola transacción el
certificado anterior pasa a \est{revocado} y se emite el nuevo, enlazados entre sí. Si las toneladas
no cambian, el certificado sigue \est{vigente} y se le anota la enmienda. Ninguno se borra. \resp{P-08} &
Estado &
SIS: certificados \\

C5 & Solo lo suyo &
$s$ tiene una relación vigente con el titular de la finca a la que pertenece el objeto: es el titular
o un representante que este autorizó. El titular puede ser una persona física o una sociedad, y tener
varias fincas. En mediciones, además, solo las que están \est{aprobadas}. \resp{P-15} \supu{P-16} &
Atributo &
DUE: todas sus lecturas \\

C6 & Ventana de auditoría &
la fecha actual está dentro de la ventana de vigencia de la cuenta de auditor, que GER fija al
activarla. Fuera de la ventana la cuenta no puede iniciar sesión, ni en sitio ni a distancia.
Prorrogarla es una autorización nueva de GER. \resp{P-18} &
Estado (tiempo) &
AUD: todas \\

C7 & Solicitar enmienda &
$m$ está \est{aprobada}, o su período está \est{cerrado} y $m$ quedó \est{devuelta} o
\est{excluida}. La solicita el técnico responsable de $m$ (TEC), un verificador (VER) o la Gerencia
(GER), con el error que encontró y su evidencia. \resp{P-07} &
Estado e historial &
TEC, VER, GER: enmiendas \\

C8 & Autorización final &
$e$ tiene dictamen técnico favorable (C10) y $s$ no la solicitó, no la revisó ni figura en
$\mathit{Autores}(m)$. Si $e$ revoca un certificado, $s$ deja registrado si ya fue comercializado y qué
se acordó con el dueño. \resp{P-07, P-08} &
Historial y estado &
GER: enmiendas \\

C9 & Corregir dentro del plazo &
$e$ está \est{solicitada}, $s$ es el técnico responsable de $m$ o el que SUP designó en su ausencia,
y no han pasado 10 días hábiles desde la solicitud. La corrección se guarda como versión
\emph{propuesta}: la versión vigente no cambia hasta la autorización final. \resp{P-03} &
Estado (tiempo) &
TEC: enmiendas \\

C10 & Revisión técnica independiente &
$e$ tiene su corrección registrada, $s \notin \mathit{Autores}(m)$ y $s$ no es quien solicitó $e$. Si
la enmienda es \emph{material} (cambia un dato que entra en el cálculo de las toneladas, o $m$ ya
tiene un certificado emitido), $s$ tampoco es quien aprobó la versión que se corrige. La
clasificación la calcula el sistema al comparar las dos versiones. \resp{P-07} &
Historial &
VER: enmiendas \\

C11 & Catálogo con segunda firma &
el alta de una finca, el cambio de titular y todo cambio de geometría o área que propone ADM quedan en
estado \est{propuesto} y solo surten efecto cuando GER los aprueba. El sistema rechaza una parcela
cuya geometría se traslape con otra parcela activa. Una parcela con mediciones o certificados no se
edita: se crea una versión nueva. \supu{P-21} &
Estado &
ADM y GER: fincas \\

C12 & Nadie se administra a sí mismo &
la cuenta afectada no es la de $s$. Asignar los roles TEC, VER, SUP, GER, ADM o AUD es una propuesta de
ADM que solo surte efecto cuando la aprueba GER, y quien aprueba no es ni quien propone ni el titular de
la cuenta. La habilitación temporal de técnico para un verificador sigue la misma vía y lleva fecha
de fin. El sistema rechaza toda asignación que viole las restricciones de la
sección~\ref{sec:exclusion}. \resp{P-01} \supu{P-20} &
Atributo e historial &
ADM y GER: cuentas y roles \\

C13 & Credenciales y declaración propias &
la cuenta es la de $s$ y el cambio se limita a su contraseña, sus factores de autenticación y su
declaración de vínculos con titulares. La declaración se puede ampliar; retirar un vínculo lo aprueba
GER. &
Atributo &
Todos los roles: cuentas \\

C14 & Apertura y cierre entre dos &
ADM registró la apertura o el cierre con las fechas que indicó la coordinación de certificación, y
GER lo aprueba: ninguno de los dos pasos basta solo. No existe la operación \emph{reabrir}. Al cerrar,
lo \est{enviado} sigue su revisión (C3), los borradores que nunca se enviaron pasan a \est{excluida}
y lo \est{devuelto} solo cambia por enmienda. \resp{P-05, P-06} &
Estado &
ADM y GER: períodos \\

C15 & Alertas ajenas &
el sujeto de la alerta no es $s$. SUP atiende las de los técnicos que tiene a cargo; GER, las de todos
los demás y las que SUP no resuelva en dos días hábiles. Las alertas sobre GER se notifican a la
Gerencia General y quedan visibles para el auditor externo. \resp{P-10} &
Atributo &
SUP y GER: alertas \\

C16 & Supresión condicionada &
no hay obligación vigente de conservar el dato: los datos personales ligados a un certificado emitido
se conservan mientras dure la trazabilidad exigida y después se desasocian del titular. &
Estado (tiempo) &
ADM: datos personales \\

C17 & Abstención por vínculo &
$s$ no tiene declarado un vínculo familiar o comercial con el titular de la finca del objeto. Quien
lo tiene no mide, no aprueba, no asigna, no revisa ni autoriza nada sobre esa finca: el sistema lo
pasa a otra persona. \resp{P-15} &
Atributo &
Las que indica la leyenda de la matriz \\

C18 & Exportar evidencia &
se cumple C6 y lo exportado es un paquete que genera el sistema: va sellado, lleva la identidad del
auditor y su huella queda en la bitácora. El titular aparece por su código; los datos personales no
se exportan. \resp{P-18} &
Estado (tiempo) &
AUD: todas menos datos personales y cuentas \\

C19 & Asignar con vigencia &
el asignado tiene el rol TEC vigente, propio o por habilitación temporal, y no es $s$. La asignación
es ordinaria, por el período, o temporal, con motivo y fecha de fin: 15 días como máximo, renovable.
\resp{P-05, P-11} &
Atributo y tiempo &
SUP: asignaciones \\

C20 & Comprobación pública &
quien consulta presenta el código impreso en el certificado. Recibe solo si el sello es válido, el
estado, las toneladas, el período y, si fue revocado, el código del que lo reemplaza. No hay
listados, búsquedas ni datos del titular. \resp{P-04} &
Atributo &
PUB: certificados \\
\end{longtable}
}

\textbf{Controles de historial sobre las lecturas marcadas con \cnd{H}.} Su definición completa está en el
documento de control de acceso; aquí se resume lo necesario para comprobarlos. Los umbrales son
parámetros por rol que aprueba GER, no constantes. \resp{P-02}

{\small
\renewcommand{\arraystretch}{1.25}
\begin{tabularx}{\textwidth}{@{}L{0.75cm} L{2.0cm} Y L{4.6cm}@{}}
\toprule
\cab{Id.} & \cab{Tipo} & \cab{Regla} & \cab{Decisión} \\
\midrule
H1 & Frecuencia &
Número de fincas distintas cuyo detalle consultó $s$ en las últimas 24 horas, comparado con los tres
umbrales de su rol. Los intentos denegados por C1 o C5 cuentan doble. \resp{P-10} \supu{P-13} &
Conceder; conceder y avisar; conceder con segundo factor y motivo; o restringir las fincas nuevas
hasta el siguiente día hábil. \\
H2 & Tiempo &
Una acción de decisión se pide fuera de la jornada del rol. \supu{P-12} &
Conceder con segundo factor y motivo. \\
H3 & Combinada &
Fuera de jornada, más fincas distintas en 60 minutos que el umbral (inicial: 10) y la mayoría nunca
consultadas por $s$ en los últimos 90 días. \resp{P-10} &
Conceder con segundo factor y motivo, y avisar de inmediato. Si la pauta sigue, restringir las
fincas nuevas hasta el siguiente día hábil. \\
H4 & Secuencia &
Quien participó en un paso anterior del flujo no decide el siguiente. Es el nombre que reciben, como
control de historial, las reglas C3, C8 y C10. &
Denegar. \\
\bottomrule
\end{tabularx}
}

%=====================================================================
\section{Restricciones entre roles}
\label{sec:exclusion}
%=====================================================================

La matriz supone que una persona tiene una sola identidad en el sistema y que los roles no se
acumulan. Son las restricciones del nivel RBAC$_2$ de Sandhu \emph{et al.} (1996), y el sistema las
valida cada vez que se asigna un rol (regla C12).

\begin{center}
\small\rmfamily
\renewcommand{\arraystretch}{1.25}
\begin{tabular}{@{}l|c|c|c|c|c|c|c@{}}
 & TEC & VER & SUP & GER & ADM & AUD & DUE \\
\hline
TEC & --- & \textbf{D} & $\times$ & $\times$ & $\times$ & $\times$ & $\times$ \\
\hline
VER & \textbf{D} & --- & $\times$ & $\times$ & $\times$ & $\times$ & $\times$ \\
\hline
SUP & $\times$ & $\times$ & --- & $\times$ & $\times$ & $\times$ & $\times$ \\
\hline
GER & $\times$ & $\times$ & $\times$ & --- & $\times$ & $\times$ & $\times$ \\
\hline
ADM & $\times$ & $\times$ & $\times$ & $\times$ & --- & $\times$ & $\times$ \\
\hline
AUD & $\times$ & $\times$ & $\times$ & $\times$ & $\times$ & --- & $\times$ \\
\hline
DUE & $\times$ & $\times$ & $\times$ & $\times$ & $\times$ & $\times$ & --- \\
\end{tabular}

\vspace{0.3em}
{\footnotesize $\times$ = exclusión estática: ninguna persona puede tener ambos roles.\quad
\textbf{D} = exclusión dinámica: una persona puede tener ambos, pero nunca activos en la misma sesión.}
\end{center}

\begin{itemize}
\item \textbf{Un rol por persona, con una sola excepción.} El cliente confirmó que algunos verificadores
      fueron técnicos y que, en situaciones excepcionales, pueden colaborar midiendo. Un verificador
      puede recibir una \textbf{habilitación temporal de técnico}, con fecha de fin, que propone ADM y
      aprueba GER (C12). Mientras dure, entra como técnico o como verificador, nunca como ambos a la
      vez. La separación no descansa en esa casilla sino en la regla C3, que mira cada medición: lo
      que midió con la habilitación no lo podrá aprobar nunca. \resp{P-01}
\item \textbf{Una identidad por persona.} La cuenta se liga al número de identificación; el sistema no
      admite dos cuentas activas para la misma persona.
\item \textbf{Cardinalidad.} Al menos dos cuentas VER activas mientras haya mediciones sin dictamen,
      para que C3 nunca deje una medición sin nadie que pueda aprobarla. Es el mínimo que exige una
      revisión independiente, no un supuesto de tamaño: ninguna regla depende de cuántos técnicos,
      verificadores o fincas haya. El número de personas con rol GER o ADM tiene un tope que fija la
      Gerencia General. \resp{P-02}
\item \textbf{Dueños y personal interno.} DUE es excluyente con todos los roles internos. Si un empleado
      es titular de una finca certificada, el acceso como dueño lo ejerce un representante. Y como
      no existe un registro que descarte los vínculos familiares o comerciales, todo el personal
      declara los suyos al ingresar y antes de cada período; la regla C17 lo aparta de esas fincas.
      \resp{P-15}
\end{itemize}

%=====================================================================
\section{Dónde se ve la separación de responsabilidades}
%=====================================================================

La acción que ninguna persona puede completar sola es la \textbf{emisión de un certificado}. En la
matriz eso se lee siguiendo la columna \emph{Mediciones} hacia \emph{Certificados}:

\begin{center}
\small
\renewcommand{\arraystretch}{1.25}
\begin{tabularx}{\textwidth}{@{}L{3.3cm} L{4.6cm} Y@{}}
\toprule
\cab{Paso del flujo} & \cab{Celda} & \cab{Qué impide} \\
\midrule
Registrar y enviar & (TEC, Mediciones): crear, modificar, enviar\cnd{C2} &
El técnico no tiene \emph{aprobar}. Después de enviar ya no puede modificar. \\
Aprobar o devolver & (VER, Mediciones): aprobar, devolver\cnd{C3} &
El verificador no tiene \emph{crear} ni \emph{modificar}, y no aprueba una medición en la que
participó, aunque haya sido en otro puesto o con una habilitación temporal. \\
Emitir & (SIS, Certificados): emitir\cnd{C4} &
Ninguna persona tiene \emph{emitir}. El sistema emite solo si los dos pasos anteriores los hicieron
personas distintas y no hay otro certificado vigente para la misma parcela y período. \\
\bottomrule
\end{tabularx}
\end{center}

El mismo patrón, quien inicia no cierra, se repite en otras cinco columnas:

\begin{center}
\small
\renewcommand{\arraystretch}{1.25}
\begin{tabularx}{\textwidth}{@{}L{3.3cm} Y L{2.2cm}@{}}
\toprule
\cab{Columna} & \cab{Quién hace cada paso} & \cab{Reglas} \\
\midrule
Enmiendas & Solicita TEC, VER o GER; corrige TEC; revisa un VER que no es autor ni solicitante; autoriza GER, que no solicitó ni revisó. Cuatro pasos y al menos tres personas. & C7 a C10 \\
Períodos de medición & Registra ADM; abre y cierra GER. & C14 \\
Asignaciones & Asigna SUP; mide TEC. Nadie se asigna a sí mismo. & C19 \\
Fincas y parcelas & Propone ADM; aprueba GER. & C11 \\
Cuentas y roles & Propone ADM; aprueba GER; nadie toca su propia cuenta. & C12 \\
\bottomrule
\end{tabularx}
\end{center}

\textbf{Registrar, modificar, revisar y aprobar son cuatro permisos distintos.} El cliente pidió que el
diseño los distinga, porque al cerrar un período pueden quedar mediciones pendientes de revisión.
\resp{P-06} Así se comporta cada uno:

\begin{center}
\small
\renewcommand{\arraystretch}{1.25}
\begin{tabularx}{\textwidth}{@{}L{2.7cm} L{2.6cm} Y Y@{}}
\toprule
\cab{Acción} & \cab{Verbo y celda} & \cab{Con el período abierto} & \cab{Con el período cerrado} \\
\midrule
Registrar & crear (TEC) & El técnico asignado. & Nadie. \\
Modificar & modificar (TEC) & El autor, si está en borrador o devuelta. & Nadie. Solo por enmienda (C7 a C10). \\
Revisar & leer, devolver (VER) & Cualquier verificador. & Cualquier verificador, sobre lo que se envió a tiempo. \\
Aprobar & aprobar (VER) & Un verificador que no participó. & Igual: el cierre no lo impide. \\
\bottomrule
\end{tabularx}
\end{center}

%=====================================================================
\section{Justificación de mínimo privilegio}
%=====================================================================

Por cada permiso amplio de la matriz, por qué no puede ser más estrecho y qué lo acota.

{\small
\renewcommand{\arraystretch}{1.25}
\begin{longtable}{@{}L{3.9cm} L{6.3cm} L{5.4cm}@{}}
\toprule
\cab{Permiso amplio} & \cab{Por qué no puede ser más estrecho} & \cab{Qué lo acota} \\
\midrule
\endfirsthead
\toprule
\cab{Permiso amplio} & \cab{Por qué no puede ser más estrecho} & \cab{Qué lo acota} \\
\midrule
\endhead
\bottomrule
\endfoot
AUD lee los diez objetos y exporta evidencia &
El cliente lo exige: los auditores «necesitan ver absolutamente todo» y, en ocasiones, llevarse
evidencia para sus informes. &
Ningún verbo de escritura; ventana de vigencia (C6); cuentas nominales; la exportación es un paquete sellado, registrado y sin datos personales (C18). \\
GER lee todo el negocio y la bitácora &
Autoriza enmiendas, atiende alertas y responde ante el auditor: sin ver la medición, el certificado y el rastro no puede decidir. &
No crea, modifica ni aprueba mediciones; no asigna fincas; no emite; vigilancia H1 a H3; no resuelve sus propias alertas (C15). \\
VER lee mediciones, fincas y certificados de cualquier finca &
La cola de verificación abarca cualquier finca del período y el dictamen compara contra el historial de la parcela y sus certificados anteriores. &
Sin acceso a datos personales del dueño ni a asignaciones; vigilancia H1 a H3; abstención por vínculo (C17). \\
SUP lee el catálogo y crea asignaciones de cualquier finca &
Coordina al personal de campo y debe poder reasignar por vacaciones, incapacidades o visitas extraordinarias. &
De las mediciones solo ve estado y fechas; no ve datos personales; no se asigna a sí mismo (C19); las temporales vencen solas; cada asignación va al resumen semanal de GER. \\
ADM lee y mantiene todo el catálogo y los datos personales &
Es su función: administrar el catálogo, habilitar las cuentas de los dueños y atender sus solicitudes de rectificación. &
Sin acceso a mediciones, enmiendas, certificados ni bitácora; todo cambio de geometría o de titular requiere a GER (C11). \\
ADM crea y desactiva cuentas y registra períodos &
Alguien tiene que dar de alta y de baja al personal el mismo día en que entra o sale, y registrar las fechas que decide la coordinación. &
No asigna roles por sí solo ni toca su propia cuenta (C12); un período no se abre ni se cierra sin GER (C14). \\
TEC crea y modifica mediciones &
Es su trabajo, y debe poder corregir lo que el verificador le devuelve. &
Solo en fincas asignadas, solo sus propios borradores y solo con el período abierto (C1, C2). \\
DUE lee todo lo de sus fincas &
El cliente lo pide: «sus mediciones, sus certificados». &
Solo lo propio (C5) y solo mediciones aprobadas; ninguna escritura. \\
SIS emite y revoca certificados &
La emisión debe ocurrir sin que una persona tenga ese permiso. &
Solo bajo C4; ninguna persona puede invocar la emisión directamente. \\
PUB comprueba un certificado &
El cliente quiere poder demostrar ante terceros que un certificado es auténtico y cuál es su estado. &
Solo con el código del certificado; respuesta mínima, sin listados ni datos del titular (C20); límite de ritmo por origen. \\
\end{longtable}
}

%=====================================================================
\section{Comprobación celda por celda}
\label{sec:casos}
%=====================================================================

Para que la matriz sea verificable por alguien que no participó en el diseño, cada regla se acompaña
de al menos un caso de prueba con su resultado esperado. El equipo que construya el sistema debe
convertirlos en pruebas automáticas y ejecutar además una prueba por cada celda marcada con
---.

{\small
\renewcommand{\arraystretch}{1.22}
\begin{longtable}{@{}L{0.9cm} L{9.5cm} L{2.9cm} L{2.2cm}@{}}
\toprule
\cab{Caso} & \cab{Situación} & \cab{Resultado} & \cab{Regla} \\
\midrule
\endfirsthead
\toprule
\cab{Caso} & \cab{Situación} & \cab{Resultado} & \cab{Regla} \\
\midrule
\endhead
\bottomrule
\endfoot
T-01 & El técnico A, asignado a la finca F1 en un período abierto, crea una medición en F1. & Concede & C1, C2 \\
T-02 & El técnico A pide el detalle de la finca F2, que no tiene asignada. & Deniega & C1; suma a H1 \\
T-03 & El técnico A intenta modificar una medición suya que ya está \est{enviada}. & Deniega & C2 \\
T-04 & El técnico A intenta modificar un borrador suyo después del cierre del período. & Deniega & C2, C14 \\
T-05 & El verificador B aprueba una medición registrada por A. & Concede & C3 \\
T-06 & La verificadora C, que el año pasado era técnica y registró la versión 1 de una medición, intenta aprobarla. & Deniega & C3 \\
T-07 & Se aprueba una segunda medición de la misma parcela y período que ya tiene certificado vigente. & No emite & C4 \\
T-08 & El dueño D cambia en la dirección el identificador de su certificado por el de otro dueño. & Deniega\cnd{b} & C5 \\
T-09 & La auditora E inicia sesión una semana después de terminada su ventana. & Deniega & C6 \\
T-10 & La auditora E, dentro de su ventana, envía una petición de modificación sobre cualquier objeto. & Deniega & Fila AUD \\
T-11 & El técnico A solicita una enmienda sobre una medición que registró otro técnico. & Deniega & C7 \\
T-12 & La gerente G autoriza una enmienda que solicitó y corrigió A y que tiene dictamen favorable de C. & Concede & C8 \\
T-13 & El técnico A registra la corrección once días hábiles después de la solicitud. & Deniega & C9 \\
T-14 & El verificador B, que solicitó la enmienda, intenta dar el dictamen técnico. & Deniega & C10 \\
T-15 & El administrador propone una parcela que se traslapa con otra activa. & Deniega & C11 \\
T-16 & El administrador se asigna a sí mismo el rol VER. & Deniega & C12 \\
T-17 & Se cierra el período; después B aprueba una medición que A envió antes del cierre. & Concede & C3, C14 \\
T-18 & El verificador B consulta en un día más fincas distintas que el umbral de verificación de su rol. & Pide segundo factor y motivo & H1 \\
T-19 & Un sábado a las 23:00, B abre en una hora el detalle de 11 fincas que nunca había consultado. & Pide segundo factor y motivo, y avisa a GER & H3 \\
T-20 & El operador de plataforma intenta iniciar sesión en la aplicación. & Deniega & Fila OPE \\
\midrule
T-21 & Después del cierre, A intenta modificar una medición que B le devolvió. & Deniega: procede una enmienda & C2, C7 \\
T-22 & El administrador registra el cierre y GER aún no lo aprueba; A crea una medición. & Concede: el período sigue abierto & C14 \\
T-23 & Enmienda material: B, que aprobó la versión que se corrige, intenta dar el dictamen técnico. & Deniega & C10 \\
T-24 & La gerente G, que solicitó una enmienda, intenta autorizarla. & Deniega & C8 \\
T-25 & Se autoriza una enmienda que cambia las toneladas de un certificado vigente. & Revoca y emite el reemplazo en un solo paso & C4 \\
T-26 & La verificadora C, con habilitación temporal de técnica, registra una medición y después, en otra sesión, intenta aprobarla. & Deniega & C3 \\
T-27 & El supervisor S da a A una asignación temporal a F2 por tres días; al cuarto, A pide el detalle de F2. & Deniega & C1, C19 \\
T-28 & El supervisor S intenta asignarse a sí mismo una finca. & Deniega & C19 \\
T-29 & B declaró un vínculo con el titular de F5 e intenta aprobar una medición de F5. & Deniega & C17 \\
T-30 & A registra una medición con fecha de medición posterior a la de hoy, o fuera del período. & Deniega & C2 \\
T-31 & La auditora E exporta la evidencia de un certificado dentro de su ventana. & Concede: paquete sellado, sin datos personales & C18 \\
T-32 & Un comprador, sin cuenta, consulta el código de un certificado revocado. & Responde «revocado» y el código del reemplazo & C20 \\
T-33 & B, después de confirmar en T-19, abre otras diez fincas que nunca había consultado. & Restringe las fincas nuevas y avisa & H3 \\
T-34 & B aprueba una medición un domingo por la tarde. & Pide segundo factor y motivo & H2 \\
T-35 & El técnico A intenta cambiar la contraseña de otra cuenta. & Deniega & C13 \\
T-36 & El supervisor S intenta resolver una alerta sobre sí mismo. & Deniega & C15 \\
T-37 & El administrador intenta suprimir los datos de un titular con un certificado emitido hace tres años. & Deniega & C16 \\
\end{longtable}
}

{\footnotesize\cnd{b} La respuesta es «no encontrado», igual que si el certificado no existiera, para no
revelar la existencia de datos de otro dueño.\quad Los casos T-21 a T-37 se agregaron después de la
consulta al cliente; T-12, T-13, T-17, T-18 y T-19 cambiaron con ella.}

%=====================================================================
\section{Qué respondió el cliente y qué cambió en la matriz}
\label{sec:minuta}
%=====================================================================

La primera versión de la matriz se construyó con supuestos de trabajo. El cliente respondió por
escrito catorce de las veintiséis consultas del equipo; la minuta completa es el Anexo~A del informe
de políticas. Esto es lo que cada respuesta cambió aquí.

{\small
\renewcommand{\arraystretch}{1.22}
\begin{longtable}{@{}L{1.2cm} L{6.0cm} L{8.4cm}@{}}
\toprule
\cab{Consulta} & \cab{Lo que respondió el cliente} & \cab{Qué cambió en la matriz} \\
\midrule
\endfirsthead
\toprule
\cab{Consulta} & \cab{Lo que respondió el cliente} & \cab{Qué cambió en la matriz} \\
\midrule
\endhead
\bottomrule
\endfoot
P-01 & Un verificador puede medir, por excepción. Hace falta revisión independiente sin impedir que el personal asuma otras tareas. & La exclusión TEC--VER pasó de estática a dinámica, con habilitación temporal (C12). C3 no cambió. Caso T-26. \\
P-02 & No hay cifras de personal ni de fincas; el diseño no debe depender de cantidades fijas. & Los umbrales de H1 y H3 son parámetros por rol. El tope de GER y ADM dejó de ser un número fijo. \\
P-03 & El verificador devuelve con observaciones; no modifica mediciones ajenas. & Confirmado: (VER, Mediciones) no tiene \emph{crear} ni \emph{modificar}, y en la enmienda corrige el técnico (C9). \\
P-04 & Interesa demostrar la autenticidad del certificado, su estado y quién autorizó su emisión. & Fila PUB y regla C20: comprobación pública con el código del certificado. \\
P-05 & La coordinación de certificación decide los períodos y el personal administrativo los registra. Las fincas las asignan los supervisores de campo. & Rol nuevo, SUP, con la columna \emph{Asignaciones} (C19). \emph{Períodos}: ADM registra y GER abre y cierra (C14). GER dejó de crear asignaciones. \\
P-06 & Al cierre pueden quedar mediciones por revisar; el cierre impide registrar y modificar, no terminar la revisión. & C3 ya no exige el período abierto. C14 solo excluye los borradores nunca enviados. Casos T-17, T-21 y T-22. \\
P-07 & Solicita el técnico responsable, un verificador o el personal de certificación; autoriza la Gerencia con una revisión técnica; los tres deben ser independientes. & La enmienda cambió de orden: solicitar, corregir, revisar, autorizar. C7, C8 y C10 se reescribieron y GER ganó \emph{solicitar}. Casos T-12, T-23 y T-24. \\
P-08 & Si cambian las toneladas hay que evaluar el certificado; nunca desaparece ni hay dos vigentes. & C4: revocar y reemitir en una sola transacción, con enlace entre ambos. C8: registro de si fue comercializado. Caso T-25. \\
P-10 & Como mínimo detectar y avisar al supervisor o a la Gerencia; respuesta según la gravedad; no bloquear de inmediato. & H1 ganó un escalón de solo aviso; H3 pide confirmación antes de restringir. La columna \emph{Alertas} tiene a SUP (C15). Casos T-19 y T-33. \\
P-11 & El técnico trabaja con sus fincas asignadas, pero debe poder apoyar, sustituir y atender visitas extraordinarias. & C1 admite asignaciones temporales con fecha de fin, que crea SUP (C19). Caso T-27. \\
P-15 & Los dueños entran; hay personas físicas y sociedades, con varias fincas y representantes. No hay registro de vínculos con empleados. & C5 reconoce titular y representante. Regla nueva C17, abstención por vínculo declarado, y C13 para declararlo. Caso T-29. \\
P-18 & No hay una norma de auditoría definida. Auditan una vez al año, en sitio o a distancia, y a veces exportan evidencia. & Verbo \emph{exportar} en la fila AUD (C18). C6 aclara que una prórroga es una autorización nueva. Caso T-31. \\
P-22 & El registro puede hacerse después de la visita; importa la trazabilidad de las fechas de medición y de registro. & C2 valida la fecha de medición declarada contra el período, la asignación y la fecha de registro. Caso T-30. \\
P-23 & La operación quedará en personal que la empresa designe. & La fila OPE se mantiene sin permisos; su descripción dejó de suponer un proveedor. \\
\end{longtable}
}

\textbf{Supuestos que el equipo fijó.} El cliente no respondió estas consultas y pidió que el equipo
decida y justifique. La matriz usa estos valores; si cambian, cambia lo que indica la última columna.

\begin{center}
\small
\renewcommand{\arraystretch}{1.22}
\begin{tabularx}{\textwidth}{@{}L{1.2cm} L{6.6cm} Y@{}}
\toprule
\cab{Consulta} & \cab{Supuesto declarado} & \cab{Qué cambiaría} \\
\midrule
P-12 & La jornada es de lunes a viernes, de 6:00 a 19:00. Es un parámetro por rol. & El valor de H2 y de H3. \\
P-13 & Hasta calibrar con datos reales, los umbrales de H1 son 40, 60 y 80 fincas por día. & Los valores de H1. \\
P-16 & El dueño ve una medición solo cuando está aprobada. & C5 en la celda (DUE, Mediciones). \\
P-20 & La jefatura solicita la cuenta, ADM la crea y GER aprueba el rol. & C12 y la celda (GER, Cuentas y roles). \\
P-21 & Una hectárea se identifica por la geometría de su parcela. & La restricción de traslape de C11. \\
\bottomrule
\end{tabularx}
\end{center}

\section*{Referencias}
\addcontentsline{toc}{section}{Referencias}

{\small
\setlength{\parskip}{0.5em}
\everypar{\hangindent=1.5em\hangafter=1}

Lampson, B. W. (1974). Protection. \emph{ACM SIGOPS Operating Systems Review, 8}(1), 18--24.
\url{https://doi.org/10.1145/775265.775268}

Saltzer, J. H., \& Schroeder, M. D. (1975). The protection of information in computer systems.
\emph{Proceedings of the IEEE, 63}(9), 1278--1308. \url{https://doi.org/10.1109/PROC.1975.9939}

Sandhu, R. S., Coyne, E. J., Feinstein, H. L., \& Youman, C. E. (1996). Role-based access control
models. \emph{Computer, 29}(2), 38--47. \url{https://doi.org/10.1109/2.485845}

Vargas Rivera, A. (2026). \emph{Material de estudio, clase 10: Matriz de control de acceso}
[Material de clase]. Curso CE-1115 Seguridad de la Información, Instituto Tecnológico de Costa Rica.

}

\end{document}
